\subsection{函数-层对应}

\ssszh 如前所约，我们已经理解了复旗簇上的 D-模范畴
\[
\on{D-mod}(\sB(\CC)).
\]
接下来要说明，它与函数构成的向量空间
$$\Fun(\sB(\Fq))$$
在何种意义上类似。

\ssszh 回忆：对任意 $\Fq$-簇 $X$ 与特征零的域 $k$，我们给它配上
\[
\Fun(X(\Fq),k),
\]
即有理点构成的有限集 $X(\Fq)$ 上 $k$-值函数构成的向量空间。

生活中一个美妙的事实是：$\Fun(X(\Fq),k)$ 中的特殊函数，往往是 $X$ 上特殊层的「影子」。
下面把这一点说得更精确。

\ssszh\label{s:ladic} 为此，取 $k=\Qlb$，即 $\ell$-进数 $\mathbb{Q}_\ell$ 的代数闭包，其中
$\ell$ 是与 $q$ 互素的任意素数。

有了这样的 $k$，就有性质良好的「$X$ 上 $k$ 向量空间的 étale 局部系统」范畴。该范畴等价于
$X$ 的 étale 基本群在 $k$ 向量空间上的连续表示范畴。

\begin{re} 对不熟悉这些概念的读者，这里给一条帮助定位的注。设 $X$ 是光滑复代数簇，带有
局部系统 $\sL$，即解析拓扑意义下局部常值的有限维 $k$ 向量空间层。

要作的基本观察是：几乎不可能在 $X$ 的某个 Zariski 开子集上把 $\sL$ 平凡化。具体地说，
读者不妨想一想 $X=\mathbb{G}_m$、$\sL$ 是绕原点单值性为单位根的秩一局部系统的情形。

然而，若基本群的相应表示通过某个有限商分解，则拉回到 $X$ 的有限覆叠空间 $\tilde{X}$ 之后
可以把 $\sL$ 平凡化。关键在于，$\tilde{X}$ 又典范地是代数簇——我们强调这在一般情形下并非
显然。在上面的例子中，若单值性是 $n$ 次单位根，可以取覆叠
\[
\mathbb{G}_m\to\mathbb{G}_m,\quad z\mapsto z^n.
\]
覆叠映射 $\tilde{X}\to X$ 是 étale（即平坦且非分歧）映射的原型。要点是：若容许不仅在
Zariski 开集上、而且在 étale 覆叠上检验层的局部性质，就得到一套稳健的局部系统理论，它对
一般概形（不必在复数上）都有意义。
\end{re}

更一般地，若还想考虑 $X$ 的局部闭子簇上的局部系统，它们都自然地属于构造 étale 层的有界导出
范畴
$
\on{Sh}(X)^b,
$
并且事实上生成该范畴。

\ssszh 下面说明这一范畴在最简单的情形 $X$ 为一个点（即域 $\kappa$ 的谱）时是什么样子。

若 $\kappa$ 可分闭，其 $k$ 向量空间层构成的构造导出范畴就是 $k$ 向量空间的有界导出范畴，即
\[
\on{Sh}(\Sp\kappa)^b\simeq\on{Vect}_k^b.
\]
显然，这是因为 $\Sp\kappa$ 没有有趣的 étale 覆叠；这也应合读者对拓扑或复几何中「点上的
通常层」的直觉。强调 $k$ 起的是层的系数域的作用，而 $\kappa$ 提供层所依托的几何对象
$\Sp\kappa$。特别地，$\kappa$ 与 $k$ 之间不必有任何关系。

然而，对一般的域 $\kappa$，$\Sp\kappa$ 确有有趣的 étale 覆叠。例如给定域的可分扩张
$\kappa\to\kappa'$，映射
\[
\Sp\kappa'\to\Sp\kappa
\]
是 étale 的，而任何有限 étale 覆叠都是这类映射的不交并。因此，记 $\kappa_s$ 为 $\kappa$ 的
可分闭包、$H:=\on{Gal}(\kappa_s/\kappa)$ 为其 Galois 群，则 $\Sp\kappa$ 上的构造层与 $H$ 在
$k$ 向量空间上的有限维连续表示的有界导出范畴之间有典范等价，即
\[
\on{Sh}(\Sp\kappa)^b\simeq\on{Rep}_k(H)^b.
\]
特别地，étale 层范畴捕捉了域 $\kappa$ 的算术信息。通常拓扑中与 $\Sp\kappa$ 类比的对象不是点，
而是（射有限）群 $H$ 的分类空间 $\on{pt}/H$。在这一类比下，人们把
$\Sp\kappa_s\to\Sp\kappa$ 想成类似于 $\on{pt}\to\on{pt}/H$。

\label{ss:shvpt}

\ssszh 给定概形的映射 $f:X\to Y$，有推前与拉回函子\footnote{此外还有 $!$-推前与 $*$-拉回，
但我们用不到。}
\[
f_*:\on{Sh}(X)^b\to\on{Sh}(Y)^b\qquad\text{和}\qquad
f^!:\on{Sh}(Y)^b\to\on{Sh}(X)^b.
\]
这里 $f_*$ 是通常的层推前，$f^!$ 是所谓的例外逆像。

若 $f$ 是紧合映射（回忆这大致是说映射的纤维是射影簇），则 $f_*$ 左伴随于 $f^!$。特别地，
给定 $X$ 的闭点 $x$（剩余域 $\kappa_x$），有伴随对
\[
i_{x,*}:\on{Sh}(\Sp\kappa_x)^b\rightleftarrows\on{Sh}(X)^b:i_x^!.
\]
具体地说，$i_{x,*}$ 给出支撑在 $x$ 处的摩天楼层，其伴随 $i_x^!$ 把层送到它在 $x$ 处支撑的
（导出）截面。

\ssszh 仍设 $X$ 是 $\Fq$ 上的簇。现在可以描述所许诺的映射
\[
\on{FF}:\on{Sh}(X)^b\to\Fun(X(\Fq)),
\]
即 Grothendieck 的函数-层对应。对每个有理点 $x:\Sp\Fq\to X$，需要给出映射
\[
\on{Sh}(X)^b\to k,
\]
即相应函数在 $x$ 处的值。提醒读者：如第 \ref{s:ladic} 节，$k$ 记 $\ell$-进数的代数闭包
$\overline{\mathbb{Q}}_\ell$。为定义所求的求值，先过渡到 $x$ 处的 $!$-茎，即考虑复合
\[
\on{Sh}(X)^b\xrightarrow{i_x^!}\on{Sh}(\Sp\Fq)^b\simeq
\on{Rep}_k(\on{Gal}(\overline{\mathbb{F}}_q/\Fq))^b.
\]
由此回忆，$\on{Gal}(\overline{\mathbb{F}}_q/\Fq)$ 就是整数环的射有限完备化
$\widehat{\mathbb{Z}}$，其拓扑生成元是 Frobenius 自同构 $\on{Fr}$。

特别地，取 Frobenius 的迹，便得到所求的映射
\[
\on{Sh}(X)^b\xrightarrow{i^!_x}
\on{Rep}_k(\on{Gal}(\overline{\mathbb{F}}_q/\Fq))^b
\xrightarrow{\on{tr}(\on{Fr},-)} k.
\]
具体地说，它把层 $\sS$ 送到数
\[
\underset{j}{\Sigma}\ (-1)^j\cdot\tr(\on{Fr},H^j(i^!_x\sS)).
\]

\begin{re} 若读者觉得这太遥远，让我们接着第 \ref{ss:shvpt} 节的类比讲下去。把完备化
$\widehat{\mathbb{Z}}$ 换成通常的整数 $\mathbb{Z}$，注意 $\on{pt}/\mathbb{Z}$ 就是圆周
$S^1$。也就是说，可以把 $\Sp\Fq$ 想成某种圆周。

于是可以把 $n$ 维簇 $X\to\Fq$ 想成类似于 $2n+1$ 维实流形 $\pi:M\to S^1$（纤维化在圆周
之上），而其有理点则是一族嵌入的圆周，给出 $\pi$ 的截面。Frobenius 迹的类比会把 $M$ 上的
局部系统送到它绕这些圆周的单值性的迹。
\end{re}

\ssszh 函数-层对应的一个基本性质是：它把函数的 $*$-推前与层的 $*$-推前、以及相应的
$!$-拉回对应起来。对拉回这是按定义，对推前则需要 étale 上同调中的 Lefschetz 不动点公式。

此外，由于 $i^!$ 把 distinguished 三角形送到 distinguished 三角形，\footnote{回忆这只是说：
$\on{Sh}(X)^b$ 中的映射序列 $A\to B\to C$，在同构意义下来自复形的短正合列。更内蕴地说，
$A$ 是 $B\to C$ 的同伦核，等价地 $C$ 是 $A\to B$ 的同伦余核，即锥。} 我们得到 $k$ 向量
空间的映射
\[
\on{FF}:K(\on{Sh}(X)^b)\underset{\mathbb{Z}}{\otimes}k\to\Fun(X(\Fq)),
\]
其中 $K(-)$ 记 Grothendieck 群；因此可以合理地把 $\on{Sh}(X)^b$ 看作函数空间的
「范畴化」。这是个有用的观点，不过我们要略微提醒：上述箭头通常并不是同构。

\ssszh 下面描述这一对应的两个基本例子。

\begin{ex} 记 $\underline{k}$ 为 $X$ 上的常层。若 $X$ 光滑且 $n$ 维，则 $\on{FF}(\underline{k})$
是 $X(\Fq)$ 上取值 $q^n$ 的常函数。
\end{ex}

\begin{re} 由于我们用的是 $!$-茎，在任何簇 $X$ 上「提升」常函数 $1$ 的层不是常层，而是
对偶化层 $\omega_X$。按定义，它由沿投影 $\pi:X\to\Sp\mathbb{F}_q$ 拉回
$\on{Gal}(\overline{\mathbb{F}}_q/\mathbb{F}_q)$ 的平凡表示 $k$ 得到，即
$$\omega_X\simeq\pi^!(k).$$ 读者应当知道，许多作者在建立函数-层对应时改用 Verdier 对偶的
约定（即用 $*$-茎），此时与常函数 $1$ 对应的会是常层，而不是对偶化层。
\end{re}

\begin{ex} 给定有理点 $x:\Sp\Fq\to X$，考虑摩天楼层 $i_{x,*}k$，其中把 $k$ 看作平凡的
$\on{Gal}(\overline{\mathbb{F}}_q/\Fq)$-模，集中在同调次数零。则有
\[
\on{FF}(i_{x,*}k)=\delta_x,
\]
即该摩天楼层「提升」了该点的示性函数。
\end{ex}

尽管这些简单的例子希望已把基本想法传达给读者，我们还应提到：远为有趣的数与函数，例如
Gauss 和、Kloosterman 和、有限李型群的特征标，都有层论的对应物。另一个丰富的例子
（它催生了整体几何朗兰兹理论）是：函数域上非分歧 Hecke 特征函数，提升为有限域上曲线的
$G$-丛模空间上的 Hecke 特征层。

\ssszh 更多细节，读者可参阅 \cite{milnenotes}、\cite{milneetalebook}、
\cite{freitagbook}。

\subsection{Riemann--Hilbert 对应}

\ssszh 从 $\Fun(X(\Fq))$ 走到 $\on{Sh}(X)^b$ 之后，我们已能理解后者对任意基域上的 $X$
（不限于 $\Fq$）的意义。

在复数上，改用关系密切的构造解析层范畴
\[
\on{Sh}(X_{\on{an}})^b
\]
会更方便。这里 $X_{\on{an}}$ 记 $X$ 的解析化，即其闭点集 $X(\CC)$ 配上度量拓扑。具体地说，
$\on{Sh}(X_{\on{an}})^b$ 是 $k$ 向量空间层的导出范畴中，由「沿 $X$ 的某个 Zariski 局部闭
分层局部常值」的层生成的全子范畴。

\begin{re} 范畴 $\on{Sh}(X_{\on{an}})^b$ 是 étale 层相应范畴的扩大，本质上来自 Artin
比较定理。例如取 $X=\mathbb{G}_m$，étale 范畴中的秩一局部系统，其单值性必须落在代数整数
的单位群 $\overline{\mathbb{Z}}_\ell^{\times}$ 中；而在解析范畴中，单值性可以是
$\overline{\mathbb{Q}}_\ell^\times$ 的任意元素。
\end{re}

\ssszh 最后，让我们得到与 D-模的联系。先固定一个同构 $\overline{\mathbb{Q}}_\ell\simeq\CC$，
以便把 $\on{Sh}(X_{\on{an}})^b$ 与它在复向量空间系数下的对应物等同起来。

回忆 $\on{Sh}(X_{\on{an}})^b$ 包含支撑在 $X$ 的子簇上的局部系统这样的对象。通过把「在微分
方程组与其局部解之间过渡」这一操作精细化，它们可以转化到 D-模的全子范畴中，如下。

\ssszh 先考虑构造导出范畴中由 $X$ 上局部系统构成的子范畴。假定 $X$ 光滑。此时有范畴的等价
\begin{equation} \label{e:vivaflatsections}
\{\text{$X_{\on{an}}$ 上有限秩局部系统}\}\simeq
\{\text{$X_{\on{an}}$ 上的平坦全纯联络}\},
\end{equation}
其中右边记复几何意义下带全纯平坦联络 $\nabla$ 的全纯向量丛 $\sE$ 构成的范畴。

我们来回顾这一等价如何运作。从右到左，把联络 $(\sE,\nabla)$ 送到其平坦截面层
$\sE^\nabla$。要点是：对任意 $x\in X$、纤维 $\sE_x$ 中的任意向量，都在某个小的解析球内
唯一地延拓为平坦截面。这里既用到联络的平坦性，也用到允许解析解（而不只是多项式解），从而
本质上化归为常微分方程的情形。从左到右，把局部系统 $\sL$ 送到全纯丛
\[
\sL\underset{\ck}{\otimes}\OO_{X_{\on{an}}},
\]
其联络由 $\OO_{X_{\on{an}}}$ 上的显然联络诱导。

\ssszh 为与 $X$ 上的代数 D-模建立联系，注意用
$-\otimes_{\OO_X}\OO_{X_{\on{an}}}$ 作用（即不仅允许多项式截面，也允许全纯截面）给出自然
的解析化函子
\begin{equation} \label{e:gaga}
\{\text{$X$ 上的平坦代数联络}\}\to
\{\text{$X_{\on{an}}$ 上的平坦全纯联络}\}.
\end{equation}
强调：左边考虑的是代数几何性质的对象，右边是复解析性质的对象。

若 $X$ 紧合，则函子 \eqref{e:gaga} 是等价。这是所谓的 GAGA（即 géométrie algébrique et
géométrie analytique）定理的推论。

然而，若 $X$ 非紧合，函子 \eqref{e:gaga} 不再是等价。它仍然本质满，但通常会把许多不同的
平坦代数联络送到同一个全纯联络。原因很简单，在秩一联络上就已可见——通常全纯函数远多于
多项式函数（与 $X$ 紧合的情形不同），而它们满足的微分方程会在解析化之后导致不同代数联络
之间的等同。

\begin{ex} 设 $X=\mathbb{A}^1$，坐标 $z$，考虑左代数 D-模
\[
D_X/(D_X\cdot\partial_z)\qquad\text{和}\qquad D_X/(D_X\cdot(\partial_z-1)).
\]
具体地说，两者的底 $O_X$-模都通过作用在 $1$ 在 $D_X$ 中的像上，与 $O_X$ 有典范等同，并分别
带有联络
\[
\nabla=\partial_z\qquad\text{和}\qquad\nabla=\partial_z-1.
\]
作为代数 D-模它们不同构，但其解析化同构。事实上，这来自 $e^z$ 是全纯而非多项式的函数，
且在不计常数倍的意义下是
\[
\partial_z u=u
\]
的唯一解，即第二个联络的平坦截面。
\end{ex}

补救 \eqref{e:gaga} 之多对一现象的，是下面这个非平凡的观察。对给定的平坦全纯丛
$\sE_{\on{an}}$，在映到它的种种代数联络 $\sE$ 之中，存在唯一的 $\sE$，其平坦截面没有
「无穷远处的本质奇点」。更精确地说，$X$ 上的代数联络 $(\sE,\nabla)$ 称为{\em 正则的}，
若对 $X$ 的一个光滑紧化 $$j:X\to\overline{X}$$（其边界
$\overline{X}\setminus X$ 是简单正规交叉除子），$j_*\sE$ 上的联络形式 $j_*\nabla$ 在边界
除子上大致只有单极点；更仔细的表述见 \cite{hottatakeuchitanisaki}。可以证明这与
$\overline{X}$ 的选取无关。

\begin{ex} 在上面的例子中，可取紧化
\[
j:\mathbb{A}^1\to\mathbb{P}^1,
\]
即添上无穷远点。在无穷远附近用坐标 $w=z^{-1}$，两个联络分别形如
\[
\nabla=\partial_w\qquad\text{和}\qquad\nabla=\partial_w+w^{-2}.
\]
特别地，第二个联络有二阶极点，因此不正则。这等价于函数 $e^z$ 在 $\infty$ 处有本质奇点。
\end{ex}

当限制到正则联络时，解析化函子
\[
\{\text{$X$ 上的正则代数联络}\}\to
\{\text{$X_{\on{an}}$ 上的平坦全纯联络}\}
\]
是等价。由此与局部系统复合所得的等价
\begin{equation} \label{e:rhconn}
\nabla:\{\text{$X$ 上的正则代数联络}\}\simeq
\{\text{$X_{\on{an}}$ 上的有限秩局部系统}\}
\end{equation}
显式地由「取平坦解析截面」给出。这称为联络的 Riemann--Hilbert 对应。

\ssszh 最后，让我们描述 \eqref{e:rhconn} 到 $X$ 上一切构造层的推广。直观上，我们现在不仅
要互换局部系统与正则联络，还要互换支撑在子簇上的那些对象。

第一个猜测或许是：构造层与某个适当的 D-模范畴之间有阿贝尔范畴的等价。毕竟上面关于联络的
对应匹配的是阿贝尔范畴。例如，人们或会期待把常层 $\CC_X$ 与正则函数 $\OO_X$ 互换，并对
$X$ 的任意闭点 $x$ 把秩一摩天楼层 $\CC_x$ 与 delta D-模 $\delta_x$ 互换。
\label{s:cannotlit}
然而，这不可能字面成立。事实上，若 $X$ 的维数 $d>0$，看阿贝尔范畴中的同态，有
\begin{align*}
\on{Hom}_{\on{Sh}(X_{\on{an}})^\heartsuit}(\CC_X,\CC_x)\simeq\CC\quad
&\text{但}\quad\on{Hom}_{\on{D-mod}(X)^\heartsuit}(\OO_X,\delta_x)\simeq 0.\\
\intertext{为修正这一点，注意过渡到导出范畴时，有}
\on{Hom}_{\on{Sh}(X_{\on{an}})}(\CC_X,\CC_x)\simeq\CC\quad
&\text{和}\quad\on{Hom}_{\on{D-mod}(X)}(\OO_X,\delta_x)\simeq\CC[-d].
\end{align*}
建议读者验证这两条断言，例如对仿射直线 $X=\mathbb{A}^1$。无论如何，我们看到所求的推广必须
涉及上同调平移，因而必须在导出范畴的层面表述。

\ssszh 有鉴于此，我们得到下面的定义。正则完整 D-模的范畴，记作
$\on{D-mod}^{\on{rh}}(X)$，是 $\on{D-mod}(X)$ 中由「光滑子簇 $Z\subset X$ 上正则代数联络的
$*$-扩张」生成的满三角子范畴。

有一个（导出）平坦解析截面的函子，定义在全部 D-模的范畴上
\begin{equation} \label{e:fllat}
\nabla:\on{D-mod}(X)\to\{\text{$\CC$ 向量空间的层}\},
\end{equation}
我们稍后会细说。这里的主要结果是：\eqref{e:fllat} 限制到正则完整 D-模后，通过构造导出范畴
分解，并给出等价
\begin{equation}
\nabla:\on{D-mod}^{\on{rh}}(X)\simeq\on{Sh}(X_{\on{an}})^b.
\end{equation}
这称为 D-模的 Riemann--Hilbert 对应。

\ssszh 内蕴地，\eqref{e:fllat} 定义如下。全纯典范丛 $\omega_{X_{\on{an}}}$ 自然地是右
D-模，\eqref{e:fllat} 由与之作张量得到：
\[
\nabla:=\omega_{X_{\on{an}}}\underset{\sD_X}{\otimes}-.
\]
与李代数同调的情形一样，$\omega_{X_{\on{an}}}$ 的一个典范射影分解给出下面的表示。记
$\Omega^i_{\on{an}}$（$1\leqslant i\leqslant d$）为 $X$ 上全纯 $i$-形式层。特别地，注意
$$\Omega^d_{\on{an}}=\omega_{X_{\on{an}}}.$$ 于是 \eqref{e:fllat} 把集中在同调次数零的
D-模 $\sM$ 送到层的 de Rham 复形
\begin{equation} \label{e:flatsec}
\sM\to\sM\underset{\OO_X}{\otimes}\Omega^1_{\on{an}}\to\cdots\to
\sM\underset{\OO_X}{\otimes}\Omega^d_{\on{an}},
\end{equation}
其中 $\sM$ 放在同调次数 $-d$。具体地说，对 $\sM$ 的局部截面 $m$ 与 $X$ 上的局部坐标
$z_1,\ldots,z_d$，第一个微分把
\[
m\mapsto\underset{i}{\Sigma}\ (\partial_i\cdot m)\otimes dz_i,
\]
更高阶的微分类似。D-模的复形则送到相应 de Rham 双复形的全化。这一构造显然尊重拟同构，
并给出所求的函子 \eqref{e:fllat}。

\ssszh 下面以一些基本例子与性质结束这段讨论。

\begin{ex} Riemann--Hilbert 对应把结构层送到常层的一个平移，即
\begin{align*}
\nabla(\OO_X)&\simeq\CC_X[d_X].\\
\intertext{它把 delta D-模 $\delta_x$ 送到相应的摩天楼层，即}
\nabla(\delta_x)&\simeq\CC_x.\\
\intertext{读者不妨直接验证，并特别注意它们与第 \ref{s:cannotlit} 节讨论的相容性。}
\intertext{\hspace{1em}更一般地，对维数为 $d_Z$ 的光滑子簇 $i:Z\to X$ 与 $Z$ 上以 $\sL$
为平坦解析截面的正则联络 $\sE$，有}
\nabla(i_*\sE)&\simeq i_*\sL[d_Z].
\end{align*}
\label{ex:exrhs}
\end{ex}

\ssszh 也可以处理奇异簇 $X$ 上的 D-模。对任何嵌入 $X\hookrightarrow S$（$S$ 光滑），该范畴
典范地等同于 $S$ 上 D-模中集合论地支撑在 $X$ 上的全子范畴。Riemann--Hilbert 对应对奇异簇
同样成立，可由光滑情形推出。

例 \ref{ex:exrhs} 中的公式则是下述事实的特例：对任意代数簇的映射 $X\to Y$，
Riemann--Hilbert 对应互换 D-模与构造层的 $*$-推前与 $!$-拉回。

\ssszh 最后讨论阿贝尔范畴。子范畴 $\on{D-mod}^{\on{rh}}(X)$ 对 $\on{D-mod}(X)$ 上的标准
截断函子封闭，并典范地等同于相应正则完整 D-模阿贝尔范畴的导出范畴。由例 \ref{ex:exrhs}
可见，这个阿贝尔范畴并不与构造层的标准阿贝尔范畴互换，而是与反常层的阿贝尔范畴等同。

\ssszh 更多细节，读者可参阅 \cite{arkhipovnotes}、\cite{bernsteinnotes}、
\cite{hottatakeuchitanisaki}。

\vspace{.5cm}

\noindent{\bf 利益冲突声明。} 无利益冲突。
